\section{KKT条件}

\subsection{KKT条件的必要性}

现在我们来思考这样一个问题，对于一个具备强对偶性的优化问题，若$\vb{x}^{*}$和$\vb*{\lambda}^{*},\vb*{\nu}^{*}$分别是原问题和对偶问题的最优解，那我们可以得出什么结论？强对偶性给出$f_0(\vb{x}^{*})=g(\vb*{\lambda},\vb*{\nu})$
\begin{Split}[KKT条件不等式]
    f_0(\vb{x}^{*})=g(\vb*{\lambda},\vb*{\nu})
    &=\inf_{\vb{x}}\qty(f_0(\vb{x})+\Sum[i=1][m]\lambda_i^{*}f_i(\vb{x})+\Sum[i=1][p]\nu_i^{*}h_i(\vb{x}))\\
    &\leq f_0(\vb{x}^{*})+\Sum[i=1][m]\lambda_i^{*}f_i(\vb{x}^{*})+\Sum[i=1][p]\nu_i^{*}h_i(\vb{x}^{*})\\
    &\leq f_0(\vb{x}^{*})
\end{Split}

显然，这里的两个$\leq$都必须取$=$才能使$f_0(\vb{x}^{*})=g(\vb*{\lambda}+\vb*{\nu})$
\begin{itemize}
    \item 第一个等号成立，意味着$\grad_{\vb{x}}L(\vb{x},\vb*{\lambda}^{*},\vb*{\nu}^{*})|_{\vb{x}^{*}}=\vb{0}$。
    \item 第二个等号成立，意味着$\lambda_i f_i(\vb{x}^{*})=0$，对于$i=1,\cdots,m$。这个性质一般称为互补松弛性（Complementary Slackness）：若$\lambda_i^{*}>0$则$f_i(\vb{x}^{*})=0$。
\end{itemize}

除此之外，$\vb{x}^{*}$和$\vb*{\lambda}^{*},\vb*{\nu}^{*}$作为最优解至少也应当是可行解。这四个条件就是著名的KKT条件！

\begin{BoxDefinition}[KKT条件]
    KKT条件（Karush-Kuhn-Tucker）是指
    \begin{itemize}
        \item 原始可行性：$f_i(\vb{x})\leq 0,\ i=1,\cdots,m,\ h_i(\vb{x})=0,\ i=1,\cdots,p$
        \item 对偶可行性：$\lambda_i\geq 0,\ i=1,\cdots,m$
        \item 互补松弛性：$\lambda_if_i(\vb{x})=0,\ i=1,\cdots,m$
        \item 稳定性：$\grad_{\vb{x}}L(\vb{x},\vb*{\lambda}^{*},\vb*{\nu}^{*})|_{\vb{x}^{*}}=\vb{0}$
    \end{itemize}
\end{BoxDefinition}

\subsection{KKT条件的充分性}

\begin{BoxTheorem}[KKT条件]
    对于满足Slater条件的凸问题，$\vb{x}^{*},\vb*{\lambda},\vb*{\nu}$是最优解当且仅当其满足KKT条件。
\end{BoxTheorem}

\begin{Proof}[\cref{thm:KKT条件}]
    必要性在上一小节已经证明了。充分性的关键在于稳定性$\grad_{\vb{x}}L(\vb{x},\vb*{\lambda}^{*},\vb*{\nu}^{*})|_{\vb{x}^{*}}=\vb{0}$为何可以导出 $L(\vb{x}^{*},\vb*{\lambda}^{*},\vb*{\nu}^{*})=\inf_{\vb{x}}L(\vb{x},\vb*{\lambda}^{*},\vb*{\nu}^{*})$。而对于凸问题，梯度为零的局部最优就是全局最优！
\end{Proof}
